\HMTNarrativeHeading{De la structure déterminantale à l'orientation}

Les constantes de Planck sont présentées après les objets qui
déterminent leurs sections d’action. Le registre dodécaphasique contient le
bloc local de Hadamard ; son quotient entier a un ordre calculable.
Le lecteur déterminantal compose cette donnée avec les coordonnées déjà engendrées
et fixe une valuation décimale. L’échelle et le coefficient d’une section
remplissent des fonctions différentes, que la démonstration maintient explicites.

Le retour distingue les sections antérieure et postérieure au tour.
Ses coefficients conservent la provenance du caractère et de la
continuation régionale. L’expression de l’action en unités physiques
est introduite au moyen du système de coordonnées dimensionnelles correspondant. La
construction transmet ainsi une section, son coefficient, son échelle et la
règle de lecture. L’utilisation ultérieure de la constante conserve ces
relations.

L’angle et le contre-angle sont développés ci-dessous comme coordonnées
d’orientation. La première lecture décale les blocs d’alpha dans la base
de complétion ; la seconde compose le coefficient adimensionnel de
retour avec alpha. Le système de coordonnées circulaire exprime le résultat en degrés ou en
radians et conserve le nombre de tours lorsque le
relèvement est utilisé. Cette séquence explique pourquoi les coordonnées angulaires
sont des intermédiaires importants entre l’action et les lecteurs
géométriques ultérieurs.

Les canaux conjugués reçoivent cette orientation et évaluent, avec
l’incidence marquée, la fonctionnelle de Barbero--Immirzi.
La réalisation électronique utilise sa trajectoire et ses lecteurs de
retour, et reçoit le quotient des sections d’action. L’exposé conserve
l’ordre de ses antécédents et distingue le coefficient adimensionnel
de la grandeur d’action qu’il contribue à représenter. Une même
constante peut agir à plusieurs étapes sans perdre l’application par laquelle
elle a été obtenue.

Le parcours prépare une composition physique concrète. L'action et l'horloge déterminent une énergie ; la transduction thermique relie cette énergie à l'information ; le lecteur d'aire incorpore l'incidence marquée. La composition longitudinale conserve l'opérateur d'inversion, sa mémoire et les incidences inscrites, et détermine une longueur sur la référence radiale déclarée. Son module d'arrivée relie l'énergie, la longueur conjuguée et le rayon réduit ; les deux identités normalisées fixent le couplage commun. L'article conserve les calculs de ces compositions pour montrer comment les résultats arithmétiques et géométriques participent à une lecture physique commune.
